\SourceLeaf{410}%
\par
This can also be expressed by the following six equations
\begingroup\footnotesize
\[
\begin{aligned}
\frac{\sin A}{\sin a}=\frac{\sin B}{\sin b}=\frac{\sin C}{\sin c}&=-\frac{\cos\tfrac12(A+B+C)}{2\sin\tfrac12a.\sin\tfrac12b.\sin\tfrac12c}\\
=\frac{\cos\tfrac12(-A+B+C)}{2\sin\tfrac12a.\cos\tfrac12b.\cos\tfrac12c}&=\frac{\cos\tfrac12(A-B+C)}{2\cos\tfrac12a.\sin\tfrac12b.\cos\tfrac12c}=\frac{\cos\tfrac12(A+B-C)}{2\cos\tfrac12a.\cos\tfrac12b.\sin\tfrac12c}
\end{aligned}
\]
\endgroup
or by these
\[
\begin{aligned}
\sin A\sin b\sin c&=\sin B\sin a\sin c=\sin C\sin a\sin b\\
&=-4\cos\tfrac12(+A+B+C).\cos\tfrac12a.\cos\tfrac12b.\cos\tfrac12c\\
&=4\cos\tfrac12(-A+B+C).\cos\tfrac12a.\sin\tfrac12b.\sin\tfrac12c\\
&=4\cos\tfrac12(+A-B+C).\sin\tfrac12a.\cos\tfrac12b.\sin\tfrac12c\\
&=4\cos\tfrac12(+A+B-C).\sin\tfrac12a.\sin\tfrac12b.\cos\tfrac12c
\end{aligned}
\]
This quantity is six times the volume of the pyramid whose vertices are the three vertices of the spherical triangle and the centre of the sphere, with the radius of the sphere taken as \(=1\). Moreover, this quantity is
\[
=4\operatorname{cotang}r.\sin\tfrac12a.\sin\tfrac12b.\sin\tfrac12c
\]
where \(r\) denotes the spherical radius of the circle circumscribed about the triangle. It is also
\[
=4\cos r.\sin\alpha.\sin\tfrac12b.\sin\tfrac12c
\]
where [\(2\alpha,2\beta,2\gamma\) denote the angles between pairs of spherical radii drawn to the vertices \(A,B,C\), opposite the sides \(a,b,c\), or] \(\alpha,\beta,\gamma\) are the angles of the plane triangle \(ABC\), since \(2\sin\tfrac12a,\ 2\sin\tfrac12b,\ 2\sin\tfrac12c\) are its sides and \(4\sin\alpha.\sin\tfrac12b.\sin\tfrac12c\) is therefore twice its area. The same triangle can be regarded as the base of the pyramid above, of height \(\cos r\), which makes the result evident. [Elsewhere, \textsc{Gauss} derives from the sixfold equation above the formulas he used so often:]
\[
\begin{aligned}
\cos\tfrac12A.\sin\tfrac12(b-c)&=\sin\tfrac12(B-C).\sin\tfrac12a\\
\sin\tfrac12A.\sin\tfrac12(b+c)&=\cos\tfrac12(B-C).\sin\tfrac12a\\
\cos\tfrac12A.\cos\tfrac12(b-c)&=\sin\tfrac12(B+C).\cos\tfrac12a\\
\sin\tfrac12A.\cos\tfrac12(b+c)&=\cos\tfrac12(B+C).\cos\tfrac12a
\end{aligned}
\]
